\sectionwithlogo{Общие сведения}{\leavevmode
\begin{mplibcode}
beginfig(0);

vardef clock(expr hours, minutes)=
	numeric clockdiam;
	clockdiam=4.5cm;

	color framecolor, bgcolor;
	%framecolor=red+.85green;
	framecolor=.8white;
	%bgcolor=.45green;%.6blue;
	%bgcolor=(.0625, .125, .175);
	bgcolor=.2white;

	fill fullcircle scaled 1.04clockdiam withcolor bgcolor;

	%%%% Shadows
	draw fullcircle scaled clockdiam shifted (bp*(1, -.5)) withcolor .4bgcolor withpen pencircle scaled 1bp;

	draw fullcircle scaled .83clockdiam shifted (bp*(1, -.5)) withcolor .4bgcolor withpen pencircle scaled 2bp;

	draw fullcircle scaled .8clockdiam withpen pencircle scaled .5bp withcolor framecolor;

	for i=1 upto 12:
		for j=1 upto 4:
			draw (.39*clockdiam*dir(90-30*i-j*6))-- (.41*clockdiam*dir(90-30*i-j*6))
			shifted (bp*(1, -.5)) withcolor .4bgcolor 
				withpen pencircle scaled .5bp;
		endfor
		draw (.38*clockdiam*dir(90-30*i))-- (.41*clockdiam*dir(90-30*i))
			shifted (bp*(1, -.5)) withcolor .4bgcolor 
			withpen pencircle scaled 1.5bp;
		label("\scriptsize\rmfamily"&decimal(i), .46*clockdiam*dir(90-30*i))
			shifted (.5bp*(1, -.5)) withcolor .4bgcolor;	
		
		for j=1 upto 4:
			draw (.39*clockdiam*dir(90-30*i-j*6))-- (.41*clockdiam*dir(90-30*i-j*6))
				withpen pencircle scaled .5bp withcolor framecolor;
		endfor
		draw (.38*clockdiam*dir(90-30*i))-- (.41*clockdiam*dir(90-30*i))
			withpen pencircle scaled 1.5bp withcolor framecolor;
		label("\scriptsize\rmfamily"&decimal(i), .46*clockdiam*dir(90-30*i)) withcolor framecolor;
	endfor

	%%%%

	draw fullcircle scaled clockdiam withpen pencircle scaled 1bp withcolor framecolor;
	draw fullcircle scaled .83clockdiam withpen pencircle scaled 2bp withcolor framecolor;
	draw fullcircle scaled .8clockdiam withpen pencircle scaled .5bp withcolor framecolor;

	label("\tiny\hologo{METAPOST}", .15clockdiam*down shifted (.5bp*(1, -.5))) withcolor .4bgcolor;
	label("\tiny\hologo{METAPOST}", .15clockdiam*down) withcolor .75framecolor;

	draw ((-.03*clockdiam*dir(90-30*hours))-- (.3*clockdiam*dir(90-30*hours)))
			shifted (1bp*(1, -.5))
			withpen pencircle scaled 3bp
			withcolor .4bgcolor;
	draw (-.03*clockdiam*dir(90-30*hours))-- (.3*clockdiam*dir(90-30*hours))
			withpen pencircle scaled 3bp withcolor framecolor;
	draw (-.05*clockdiam*dir(90-6*minutes))-- (.35*clockdiam*dir(90-6*minutes))
			withpen pencircle scaled 1.8bp withcolor bgcolor;
	draw ((-.05*clockdiam*dir(90-6*minutes))-- (.35*clockdiam*dir(90-6*minutes)))
			shifted (1bp*(1, -.5))
			withpen pencircle scaled 1.5bp
			withcolor .4bgcolor;
	draw (-.05*clockdiam*dir(90-6*minutes))-- (.35*clockdiam*dir(90-6*minutes))
			withpen pencircle scaled 1.5bp withcolor framecolor;

	fill fullcircle scaled 1bp withcolor .4bgcolor;
enddef;

numeric minutes, hours;
minutes=time;
hours=minutes/60;
if hours>=12: hours:=hours-12; fi

clock(hours, minutes);

endfig;
\end{mplibcode}
}

%%%%%%%%%%%%%%%%%%%%%%%%%

\begin{frame}{Что такое \hologo{METAPOST}?}
\begin{itemize}
\item
это алгоритмический язык
\item
это система программирования (компилятор~+~библиотеки)
\end{itemize}

\hologo{METAPOST} был создал Джоном Хобби (John~D.~Hobby).

Во многом система \hologo{METAPOST} основана на системе \hologo{METAFONT},
созданной известным американским математиком Дональдом Кнутом
(Donald~E.~Knuth). \hologo{METAFONT} задуман как средство изготовления шрифтов
для системы компьютерной вёрстки \TeX.
\end{frame}

%%%%%%%%%%%%%%%%%%%%%%%%%

\begin{frame}{Возможности \hologo{METAPOST}}

\begin{itemize}
\item
создание рисунков в~векторном графическом формате PostScript

\item
сложные вычисления при подготовке рисунков

\item
широкий набор встроенных и~библиотечных команд

\item
возможность программирования пользовательских команд

\item
вставка тескта, отформатированного в~системе~\TeX

\end{itemize}
\end{frame}

%%%%%%%%%%%%%%%%%%%%%%%%%

\begin{frame}{Сфера применения \hologo{METAPOST}}

\begin{columns}

\column{.55\textwidth}
\begin{itemize}
\item
\alert<1>{подготовка чертежей и~технических рисунков}
\item
\alert<2>{диаграмм}
\item
\alert<3>{графиков функций}
\item
\alert<4>{иллюстраций по геометрии (с~формулами)}
\item
\alert<5>{логотипов}
\end{itemize}

\column{.45\textwidth}
\only<1>{\leavevmode
\begin{mplibcode}
beginfig(0);

vardef clock(expr hours, minutes)=
	numeric clockdiam;
	clockdiam=4.5cm;

	color framecolor, bgcolor;
	%framecolor=red+.85green;
	framecolor=.8white;
	%bgcolor=.45green;%.6blue;
	%bgcolor=(.0625, .125, .175);
	bgcolor=.2white;

	fill fullcircle scaled 1.04clockdiam withcolor bgcolor;

	%%%% Shadows
	draw fullcircle scaled clockdiam shifted (bp*(1, -.5)) withcolor .4bgcolor withpen pencircle scaled 1bp;

	draw fullcircle scaled .83clockdiam shifted (bp*(1, -.5)) withcolor .4bgcolor withpen pencircle scaled 2bp;

	draw fullcircle scaled .8clockdiam withpen pencircle scaled .5bp withcolor framecolor;

	for i=1 upto 12:
		for j=1 upto 4:
			draw (.39*clockdiam*dir(90-30*i-j*6))-- (.41*clockdiam*dir(90-30*i-j*6))
			shifted (bp*(1, -.5)) withcolor .4bgcolor 
				withpen pencircle scaled .5bp;
		endfor
		draw (.38*clockdiam*dir(90-30*i))-- (.41*clockdiam*dir(90-30*i))
			shifted (bp*(1, -.5)) withcolor .4bgcolor 
			withpen pencircle scaled 1.5bp;
		label("\scriptsize\rmfamily"&decimal(i), .46*clockdiam*dir(90-30*i))
			shifted (.5bp*(1, -.5)) withcolor .4bgcolor;	
		
		for j=1 upto 4:
			draw (.39*clockdiam*dir(90-30*i-j*6))-- (.41*clockdiam*dir(90-30*i-j*6))
				withpen pencircle scaled .5bp withcolor framecolor;
		endfor
		draw (.38*clockdiam*dir(90-30*i))-- (.41*clockdiam*dir(90-30*i))
			withpen pencircle scaled 1.5bp withcolor framecolor;
		label("\scriptsize\rmfamily"&decimal(i), .46*clockdiam*dir(90-30*i)) withcolor framecolor;
	endfor

	%%%%

	draw fullcircle scaled clockdiam withpen pencircle scaled 1bp withcolor framecolor;
	draw fullcircle scaled .83clockdiam withpen pencircle scaled 2bp withcolor framecolor;
	draw fullcircle scaled .8clockdiam withpen pencircle scaled .5bp withcolor framecolor;

	label("\tiny\hologo{METAPOST}", .15clockdiam*down shifted (.5bp*(1, -.5))) withcolor .4bgcolor;
	label("\tiny\hologo{METAPOST}", .15clockdiam*down) withcolor .75framecolor;

	draw ((-.03*clockdiam*dir(90-30*hours))-- (.3*clockdiam*dir(90-30*hours)))
			shifted (1bp*(1, -.5))
			withpen pencircle scaled 3bp
			withcolor .4bgcolor;
	draw (-.03*clockdiam*dir(90-30*hours))-- (.3*clockdiam*dir(90-30*hours))
			withpen pencircle scaled 3bp withcolor framecolor;
	draw (-.05*clockdiam*dir(90-6*minutes))-- (.35*clockdiam*dir(90-6*minutes))
			withpen pencircle scaled 1.8bp withcolor bgcolor;
	draw ((-.05*clockdiam*dir(90-6*minutes))-- (.35*clockdiam*dir(90-6*minutes)))
			shifted (1bp*(1, -.5))
			withpen pencircle scaled 1.5bp
			withcolor .4bgcolor;
	draw (-.05*clockdiam*dir(90-6*minutes))-- (.35*clockdiam*dir(90-6*minutes))
			withpen pencircle scaled 1.5bp withcolor framecolor;

	fill fullcircle scaled 1bp withcolor .4bgcolor;
enddef;

numeric minutes, hours;
minutes=time;
hours=minutes/60;
if hours>=12: hours:=hours-12; fi

clock(hours, minutes);

endfig;
\end{mplibcode}
}
\only<2>{\leavevmode
\begin{mplibcode}
input piechartmp;

beginfig(0)

Segment(30, "A");
Segment(35, "B");
Segment(50, "C");
Segment(150, "D");
PieChart(2cm, .3, 60, 180, .1);

endfig;
\end{mplibcode}
}
\only<3>{\leavevmode
\begin{mplibcode}
input graphs;

beginfig(0);

drawaxes(-cm, -3cm, 3cm, 3cm);

vardef f(expr t)=
	t*(t-1)*(t-2)
enddef;

draw graph.f(-.6, 2.6, .025, cm, cm) withpen boldpen;

label.bot("$x$", (3cm, 0));
label.lft("$y$", (0,  3cm));

endfig;
\end{mplibcode}
}
\only<4>{\leavevmode
\begin{mplibcode}
beginfig(0)

z1=origin;
z2=(2cm, 0);
z3=(1.5cm, 5cm);
z4=.45[z1, z2]; % perpbase(z3, z1, z2);
z5=.15[z1, z3];
z6=.1[z2, z3];
z7=z4 reflectedabout(z1, z3);
z8=z4 reflectedabout(z2, z3);

draw z1--z2--z3--cycle withpen boldpen;
draw z3--z4 withpen thinpen;
draw z4--z5--z6--cycle withpen thinpen;
draw z3--z7--z5 withpen thinpen;
draw z3--z8--z6 withpen thinpen;
draw z7--z8 withpen thinpen;

markpath.ticks(z3--z4, 1);
markpath.ticks(z3--z7, 1);
markpath.ticks(z3--z8, 1);
markpath.ticks(z5--z4, 2);
markpath.ticks(z5--z7, 2);
markpath.ticks(z6--z4, 3);
markpath.ticks(z6--z8, 3);

markangle.arc.ccw(z3, z1, z4, 1);
markangle.arc.ccw(z3, z7, z1, 1) scale=>1.05; % TODO extra=>1;
markangle.arc.ccw(z3, z4, z2, 2); % TODO extra=>1;
markangle.arc.ccw(z3, z2, z8, 2) scale=>1.05;

label.bot("$A$", z1);
label.bot("$B$", z2);
label.top("$C$", z3);
label.bot("$R$", z4);
label.ulft("$Q$", z5);
label.urt("$P$", z6);
label.bot("$R'$", z7);
label.bot("$R''$", z8);

endfig
\end{mplibcode}
}
\only<5>{\leavevmode
\begin{mplibcode}
beginfig(0)

numeric w, circ[], h, u;
u:=3bp;

w:=1200/36u;
circ0=.4u;
circ1=1.2u;
h:=w;

color gridcolor, integralcolor, moebiuscolor;
gridcolor=(.125, .25, .35);
moebiuscolor=1.5gridcolor;
integralcolor=1.5gridcolor;

wh[0]:=w; wh[1]:=h;
for i=0 upto 1:
  k[i]:=wh[i]/35; p[i]:=floor(4k[i]);
  s[i]:=floor(p[i]/2.5);
  q[i]:=wh[i]-7p[i]-2s[i]; r[i]:=round(q[i]/8); wh[i]:=2s[i]+7p[i]+8r[i];
endfor
w:=wh0; h:=wh1;
%TZh:=h/aspect_ratio;

path P;
%TZ currentpicture:=unitpixel xscaled wh0 yscaled wh1;
%TZ currentpicture:=unitsquare xscaled wh0 yscaled wh1;
%TZ P:=unitpixel xscaled (7p0+8r0) yscaled (7p1+8r1) shifted (s0,s1);
P:=unitsquare xscaled (7p0+8r0) yscaled (7p1+8r1) shifted (s0,s1);
%TZ addto currentpicture also -P;
fill P withcolor gridcolor;

P:=unitsquare xscaled p0 yscaled p1;
for m=0 upto 6:
	for n=0 upto 6:
		fill (P shifted ((s0+r0,s1+r1)+(m*(p0+r0),n*(p1+r1))))
			withcolor background;
	endfor
endfor

%%   Möbius band coordinates

z1=(0.11881w, .24587h);  z2=(.29938w, .59362h);  z3=(.56688w, .88119h);
z4=(0.72069w, .85444h);  z5=(.45319w, .56688h);  z6=(.28600w, .27931h);
z7=(0.12550w, .21912h);  z8=(.28600w, .24587h);  z9=(.58025w, .33281h);
z10=(0.88788w, .46656h); z11=(.83438w, .29938h); z12=(.54681w, .17900h);
z13=(0.21912w, .09206h); z14=(.76081w, .42644h); z15=(.64044w, .74075h);
z16=(0.74075w, .83438h); z17=(.88119w, .49331h); z18=(.82769w, .67388h);
z19=(0.70731w, .59362h); z20=(.60031w, .74744h); z21=(.76750w, .39969h);

%%   Integeral sign coordinates

wide:=.1w; bold:=.07w; vair:=.05h; bulb:=.07h;
penpos31(bulb,-90); penpos37(bulb,-90);
penpos32(vair,-90); penpos36(vair,-90);
penpos33(bold,0); penpos35(bold,0);
penpos34(wide,0);
z34=(.5w,.5h); z32=(.5w,1.03h);
z33+z35=z32+z36=z31+z37=(w,h);
y33=.8h; x33l=x34l-0.03w;
x31=.55w; y31l=y32l;

path S;
S=z34l..tension 4..z35l...z36l...z37l...z37r...z36r..tension 1.1..z35r...z34r..
  tension 4..z33r...z32r...z31r...z31l...z32l..tension 1.1..z33l...cycle;

fill z1--z6--z8--z7--cycle withcolor background;
fill z4--z16--z15--z20--cycle withcolor background;
fill z10--z17--z14--z21--cycle withcolor background;
%erase fill flex(z1,z2,z3) & z3--z4 & flex(z4,z5,z6) & z6--z1 & cycle;
pickup pencircle scaled circ0;
fill flex(z1,z2,z3) & z3--z4 & flex(z4,z5,z6) & z6--z1 & cycle
		withcolor moebiuscolor;
draw flex(z1,z2,z3) & z3--z4 & flex(z4,z5,z6) & z6--z1 & cycle
		withcolor black;
%erase fill flex(z14,z19,z15) & z15--z16 & flex(z16,z18,z17) & z17--z14 & cycle;
fill flex(z14,z19,z15) & z15--z16 & flex(z16,z18,z17) & z17--z14 & cycle
		withcolor moebiuscolor;
draw flex(z14,z19,z15) & z15--z16 & flex(z16,z18,z17) & z17--z14 & cycle
		withcolor black;
pickup pencircle scaled circ1;
fill S withcolor integralcolor;
draw S withcolor black;
pickup pencircle scaled circ0;
filldraw S withcolor integralcolor;
%erase fill flex(z7,z9,z10) & z10--z11 & flex(z11,z12,z13) & z13--z7 & cycle;
fill flex(z7,z9,z10) & z10--z11 & flex(z11,z12,z13) & z13--z7 & cycle
		withcolor moebiuscolor;
%pickup pencircle scaled circ1;
draw flex(z7,z9,z10) & z10--z11 & flex(z11,z12,z13) & z13--z7 & cycle
		withcolor black;

endfig
\end{mplibcode}
}

\end{columns}

\transdissolve<2->

\end{frame}

%%%%%%%%%%%%%%%%%%%%%%%%%

\begin{frame}{Особенности \hologo{METAPOST}}

\hologo{METAPOST} \alert{не} является системой WYSIWYG (What You See Is What
You Get~— «вы получаете то, что видите»). Это значит, что рисунок готовится как
программа, которая компилируется, и~в~результате компиляции получается файл
в~формате PostScript, содержащий изображение.

\end{frame}

%%%%%%%%%%%%%%%%%%%%%%%%%

\begin{frame}{PostScript}

PostScript~— это не только графический формат, но и~алгоритмический язык. В~нём
заложены богатые графические и~вычислительные возможности.

Однако из-за особенностей этого языка на нём трудно программировать «вручную».
Такое программирование требует от программиста совершенно иного мышления,
нежели программирование на «традиционных» алгоритмических языках.

Чаще всего программы на PostScript генерируются автоматически другими
программами.

\end{frame}

%%%%%%%%%%%%%%%%%%%%%%%%%

\begin{frame}{Предназначение PostScript}
Главное предназначение PostScript~— управление растровыми устройствами, такими
как принтеры, экран монитора или графический файл.

Некоторые устройства вроде дорогих профессиональных принтеров имеют встроенный
интерпретатор PostScript. Другие понимают более примитивный управляющий язык,
поэтому для управления ими требуется программа"=переводчик (транслятор).
Наиболее популярный транслятор языка PostScript~— программа GhostScript.
\end{frame}

%%%%%%%%%%%%%%%%%%%%%%%%%

\begin{frame}{Пример на языке PostScript}
\begin{columns}
\column{.3\textwidth}
\includegraphics{pstest}
\column{.7\textwidth}
\begin{programlisting}
\%!PS-Adobe\par
\leavevmode\par
5 setlinewidth\par
1 .388235294117647 .27843137254902\par
~~~~setrgbcolor\par
\leavevmode\par
newpath\par
8 dup moveto\par
72 dup rlineto\par
-72 0 rlineto\par
72 -72 rlineto\par
stroke\par
\leavevmode\par
showpage
\end{programlisting}
\end{columns}

\end{frame}

%%%%%%%%%%%%%%%%%%%%%%%%%

\begin{frame}{Где взять \hologo{METAPOST}}
\begin{itemize}
\item
\hologo{METAPOST} для Linux, Windows, Mac~OS~X входит в~пакет программ
\TeX~Live:
\url{https://tug.org/texlive/}
\item
\hologo{METAPOST} для Windows входит в~пакет программ MiK\TeX:
\url{https://miktex.org/}
\end{itemize}
\end{frame}

\begin{frame}
\frametitle{Что почитать про \hologo{METAPOST}\only<2->{
	{\mdseries\itshape(продолжение)}}}
\begin{thebibliography}{1}
\small
\only<1>{
\bibitem{bib:Hobby}
A User’s Manual for MetaPost
\newblock
Hobby, John~D.
\newblock
\url{https://www.tug.org/docs/metapost/mpman.pdf}

\bibitem{bib:Hobby}
\hologo{METAPOST}. Руководство пользователя
\newblock
Hobby, John D.
\newblock
\url{http://www.ctan.org/tex-archive/info/metapost/doc/russian/mpman-ru/mpman-ru.pdf}

\bibitem{bib:Grogono}
\hologo{METAPOST}: A~Reference Manual
\newblock
Grogono, Peter
\newblock
\url{http://users.encs.concordia.ca/~grogono/Writings/mpref.pdf}

\bibitem{bib:Wiki}
MetaPost — Википедия
\newblock
\url{http://ru.wikipedia.org/wiki/MetaPost}
\newblock
Статья в русской Википедии

\bibitem{bib:Zoonekynd}
Exemples d’utilisation de Métapost
\newblock
Zoonekynd, Vincent
\newblock
\url{http://zoonek.free.fr/LaTeX/Metapost/metapost.html}

}

\only<2>{%
\bibitem{bib:Baldin}
Создание иллюстраций в~\hologo{METAPOST}
\newblock
Балдин, Е. М.
\newblock
\url{http://www.inp.nsk.su/~baldin/mpost/}
\newblock
Статья для журнала «Linux Format»

\bibitem{bib:Heck}
Learning \hologo{METAPOST} by Doing
\newblock
Heck, André
\newblock
\url{http://remote.science.uva.nl/~heck/Courses/mptut.pdf}

\bibitem{bib:Urs}
MetaPost: A Very Brief Tutorial
\newblock
Urs, Oswald
\newblock
\url{http://www.ursoswald.ch/metapost/tutorial.pdf}

\bibitem{bib:Hurlin}
Practical introduction to~\hologo{METAPOST}
\newblock
Hurlin, Clément
\newblock
\url{http://www-sop.inria.fr/everest/Clement.Hurlin/misc/Practical-introduction-to-MetaPost.pdf}

\bibitem{bib:Knuth:TheMETAFONTbookWilliams}
Всё про \hologo{METAFONT}
\newblock
Кнут, Дональд Э.
\newblock
Москва~: Вильямс, 2003.~— ISBN 5-8459-0442-0 (рус.)
}
\end{thebibliography}
\end{frame}
